\RequirePackage{iftex}
\PassOptionsToPackage{table}{xcolor}
\documentclass{bmstu}

\ifPDFTeX
    \input{glyphtounicode.tex}
    \pdfgentounicode=1
\else
    \usepackage{fontspec}
    \defaultfontfeatures{Ligatures=TeX}
    \setmainfont{Times New Roman}[
        Path=/mnt/c/Windows/Fonts/,
        Extension=.ttf,
        UprightFont=times,
        BoldFont=timesbd,
        ItalicFont=timesi,
        BoldItalicFont=timesbi
    ]
    \setsansfont{Arial}[
        Path=/mnt/c/Windows/Fonts/,
        Extension=.ttf,
        UprightFont=arial,
        BoldFont=arialbd,
        ItalicFont=ariali,
        BoldItalicFont=arialbi
    ]
    \setmonofont{Courier New}[
        Path=/mnt/c/Windows/Fonts/,
        Extension=.ttf,
        UprightFont=cour,
        BoldFont=courbd,
        ItalicFont=couri,
        BoldItalicFont=courbi
    ]
\fi

\usepackage{adjustbox}
\usepackage{amsfonts}
\usepackage{amsmath}
\usepackage{array}
\usepackage{blindtext}
\usepackage{caption}
\usepackage{chngcntr}
\usepackage{csvsimple}
\usepackage{enumitem}
\usepackage{etoolbox}
\usepackage{float}
\usepackage{geometry}
\usepackage{graphicx}
\usepackage{hhline}
\usepackage{indentfirst}
\usepackage{lastpage}
\usepackage{listings}
\lstdefinelanguage{Go}{
    morekeywords={break,default,func,interface,select,case,defer,go,map,struct,chan,else,goto,package,switch,const,fallthrough,if,range,type,continue,for,import,return,var,nil,true,false},
    sensitive=true,
    morecomment=[l]{//},
    morecomment=[s]{/*}{*/},
    morestring=[b]",
}
\usepackage{microtype}
\usepackage{multirow}
\usepackage{pdfpages}
\usepackage{ragged2e}
\usepackage{rotating}
\usepackage{setspace}
\usepackage{siunitx}
\usepackage{tabularx}
\usepackage{tcolorbox}
\usepackage{titletoc}
\usepackage{totcount}
\usepackage{refcount}
\usepackage{titlesec}
\usepackage{ulem}
\usepackage{url}
\usepackage{tikz}
\usetikzlibrary{positioning, shapes, arrows.meta, fit, calc}
\usepackage{xparse}

\addto\captionsrussian{
    \renewcommand{\contentsname}{СОДЕРЖАНИЕ}
    \renewcommand{\bibname}{СПИСОК ИСПОЛЬЗОВАННЫХ ИСТОЧНИКОВ}
}

\addto\captionsrussian{
    \renewcommand{\contentsname}{СОДЕРЖАНИЕ}
    \renewcommand{\bibname}{СПИСОК ИСПОЛЬЗОВАННЫХ ИСТОЧНИКОВ}
}

% Словоформы для количественных числительных
\makeatletter
\newcommand{\ruquantity}[4]{%
    \number\numexpr#1\relax~\ruquantityword{#1}{#2}{#3}{#4}%
}
\newcommand{\ruquantityword}[4]{%
    \begingroup
    \count255=\numexpr#1\relax
    \count254=\count255
    \divide\count254 by 100
    \multiply\count254 by 100
    \advance\count255 by -\count254
    \ifnum\count255>10
        \ifnum\count255<20
            #4%
        \else
            \count255=\numexpr#1\relax
            \count254=\count255
            \divide\count254 by 10
            \multiply\count254 by 10
            \advance\count255 by -\count254
            \ifnum\count255=1 #2% 
            \else
                \ifnum\count255>1
                    \ifnum\count255<5 #3% 
                    \else #4% 
                    \fi
                \else #4% 
                \fi
            \fi
        \fi
    \else
        \count255=\numexpr#1\relax
        \count254=\count255
        \divide\count254 by 10
        \multiply\count254 by 10
        \advance\count255 by -\count254
        \ifnum\count255=1 #2% 
        \else
            \ifnum\count255>1
                \ifnum\count255<5 #3% 
                \else #4% 
                \fi
            \else #4% 
            \fi
        \fi
    \fi
    \endgroup
}
\makeatother

% Форматирование секций
\titleformat{name=\chapter,numberless}[block]{\normalfont\large\bfseries\centering}{}{0pt}{}
\titleformat{\chapter}[block]{\normalfont\large\bfseries}{\thechapter}{1em}{}
\titlespacing\chapter{\parindent}{*2}{*2}

\titleformat{name=\section,numberless}[block]{\normalfont\large\bfseries\centering}{}{0pt}{}
\titleformat{\section}[block]{\normalfont\large\bfseries}{\thesection}{1em}{}
\titlespacing\section{\parindent}{*3}{*2}

\titleformat{\subsection}[hang]{\bfseries\large}{\thesubsection}{1em}{}
\titlespacing\subsection{\parindent}{*2}{*2}

\titleformat{\subsubsection}[hang]{\bfseries\large}{\thesubsubsection}{1em}{}
\titlespacing\subsubsection{\parindent}{*2}{*2}

% Отточие в оглавлении
\titlecontents{chapter}[0pt]
  {\addvspace{1em}\bfseries}
  {\contentslabel{2.3em}}
  {}
  {\normalfont\hspace{0.5em}\titlerule*[1pc]{.}\contentspage}
  []

\titlecontents{section}[1.5em]
  {\addvspace{0.5em}}
  {\contentslabel{2.5em}}
  {}
  {\normalfont\hspace{0.5em}\titlerule*[1pc]{.}\contentspage}
  []

\titlecontents{subsection}[4em]
  {}
  {\contentslabel{3em}}
  {}
  {\normalfont\hspace{0.5em}\titlerule*[1pc]{.}\contentspage}
  []

% Общее форматирование
\onehalfspacing
\frenchspacing
\setlength\parindent{1.25cm}
\hypersetup{hidelinks,unicode=true}
\microtypesetup{expansion=false}

% Настройка списков
\setlist{
    labelsep=0.3cm,
    itemsep=0pt,
    parsep=0pt,
    topsep=0pt,
    itemindent=0pt
}

\setlist[enumerate]{
    label=\arabic*),
    leftmargin=2cm,
}
\setlist[enumerate,2]{
    label=\arabic*),
    leftmargin=*,
}

\setlist[itemize]{
    label=---,
    leftmargin=2cm,
    labelsep=0.3cm
}
\setlist[itemize,2]{
    label=---,
    leftmargin=*,
    labelsep=0.3cm
}

% Настройки листингов
\lstset{
    backgroundcolor=\color{white},
    basicstyle=\footnotesize\ttfamily\linespread{0.8},
    breakatwhitespace=true,
    breaklines=true,
    abovecaptionskip=8pt,
    belowcaptionskip=10pt,
    captionpos=b,
    commentstyle=\color{green!50!black},
    escapeinside={\#*}{*)},
    frame=single,
    inputencoding=utf8,
    keywordstyle=\color{blue},
    language=C,
    numbers=left,
    numbersep=5pt,
    numberstyle=\tiny,
    showstringspaces=false,
    stringstyle=\color{red!90!black},
    tabsize=4,
}
\lstset{morekeywords={size_t, likely}}
\lstset{
    literate=
    {а}{{\selectfont\char224}}1
    {б}{{\selectfont\char225}}1
    {в}{{\selectfont\char226}}1
    {г}{{\selectfont\char227}}1
    {д}{{\selectfont\char228}}1
    {е}{{\selectfont\char229}}1
    {ё}{{\"e}}1
    {ж}{{\selectfont\char230}}1
    {з}{{\selectfont\char231}}1
    {и}{{\selectfont\char232}}1
    {й}{{\selectfont\char233}}1
    {к}{{\selectfont\char234}}1
    {л}{{\selectfont\char235}}1
    {м}{{\selectfont\char236}}1
    {н}{{\selectfont\char237}}1
    {о}{{\selectfont\char238}}1
    {п}{{\selectfont\char239}}1
    {р}{{\selectfont\char240}}1
    {с}{{\selectfont\char241}}1
    {т}{{\selectfont\char242}}1
    {у}{{\selectfont\char243}}1
    {ф}{{\selectfont\char244}}1
    {х}{{\selectfont\char245}}1
    {ц}{{\selectfont\char246}}1
    {ч}{{\selectfont\char247}}1
    {ш}{{\selectfont\char248}}1
    {щ}{{\selectfont\char249}}1
    {ъ}{{\selectfont\char250}}1
    {ы}{{\selectfont\char251}}1
    {ь}{{\selectfont\char252}}1
    {э}{{\selectfont\char253}}1
    {ю}{{\selectfont\char254}}1
    {я}{{\selectfont\char255}}1
    {А}{{\selectfont\char192}}1
    {Б}{{\selectfont\char193}}1
    {В}{{\selectfont\char194}}1
    {Г}{{\selectfont\char195}}1
    {Д}{{\selectfont\char196}}1
    {Е}{{\selectfont\char197}}1
    {Ё}{{\"E}}1
    {Ж}{{\selectfont\char198}}1
    {З}{{\selectfont\char199}}1
    {И}{{\selectfont\char200}}1
    {Й}{{\selectfont\char201}}1
    {К}{{\selectfont\char202}}1
    {Л}{{\selectfont\char203}}1
    {М}{{\selectfont\char204}}1
    {Н}{{\selectfont\char205}}1
    {О}{{\selectfont\char206}}1
    {П}{{\selectfont\char207}}1
    {Р}{{\selectfont\char208}}1
    {С}{{\selectfont\char209}}1
    {Т}{{\selectfont\char210}}1
    {У}{{\selectfont\char211}}1
    {Ф}{{\selectfont\char212}}1
    {Х}{{\selectfont\char213}}1
    {Ц}{{\selectfont\char214}}1
    {Ч}{{\selectfont\char215}}1
    {Ш}{{\selectfont\char216}}1
    {Щ}{{\selectfont\char217}}1
    {Ъ}{{\selectfont\char218}}1
    {Ы}{{\selectfont\char219}}1
    {Ь}{{\selectfont\char220}}1
    {Э}{{\selectfont\char221}}1
    {Ю}{{\selectfont\char222}}1
    {Я}{{\selectfont\char223}}1
}

% Настройка подписей
\captionsetup[figure]{
    name={Рисунок},
    justification=centering,
    labelsep=endash,
    singlelinecheck=on
}
\captionsetup[table]{
    name={Таблица},
    justification=raggedright,
    labelsep=endash,
    singlelinecheck=off
}
\captionsetup[lstlisting]{
    justification=raggedright,
    labelsep=endash,
    singlelinecheck=false,
    skip=8pt
}

% Метки в библиографии
\makeatletter
\renewcommand\@biblabel[1]{#1.}
\makeatother

% Подчёркнутый текст
\NewDocumentCommand{\ulinetext}{O{3cm} O{c} m m}
{\underset{\mbox{\scriptsize #3}}{\uline{\makebox[#1][#2]{#4}}}}

% Секция без номера, но с записью в оглавлении
\newcommand{\numberlesssection}[1]{\phantomsection\section*{#1}\addcontentsline{toc}{section}{#1}}

% Подсекция без номера, но с записью в оглавлении
\newcommand{\numberlesssubsection}[1]{\phantomsection\subsection*{#1}\addcontentsline{toc}{subsection}{#1}}

% Центрирование в таблице, растянутой по ширине страницы
\newcolumntype{Y}{>{\centering\arraybackslash}X}

% Приложения
\newcounter{appsection}
\regtotcounter{appsection}
\newcommand{\russianappendixletter}[1]{%
    \ifcase#1\or А\or Б\or В\or Г\or Д\or Е\or Ж\or З\or И\or К\or Л\or М\or Н\or О\or П\or Р\or С\or Т\or У\or Ф\or Х\or Ц\or Ч\or Ш\or Щ\or Э\or Ю\or Я\else ?\fi%
}
\renewcommand{\theappsection}{\russianappendixletter{\value{appsection}}}

\newcommand{\appsection}{
    \newpage
    \refstepcounter{appsection}
    \phantomsection
    \addcontentsline{toc}{chapter}{ПРИЛОЖЕНИЕ~\theappsection}
    \begin{center}
        \textbf{ПРИЛОЖЕНИЕ~\theappsection}
    \end{center}

    \renewcommand{\thefigure}{\theappsection.\arabic{figure}}
    \renewcommand{\thetable}{\theappsection.\arabic{table}}
    \renewcommand{\thelstlisting}{\theappsection.\arabic{lstlisting}}
    \renewcommand{\theequation}{\theappsection.\arabic{equation}}

    \setcounter{figure}{0}
    \setcounter{table}{0}
    \setcounter{lstlisting}{0}
    \setcounter{equation}{0}
}

% Листинги с автопереносом
\tcbuselibrary{listings, breakable, skins}
\newtcblisting[use counter=lstlisting]{listing}[3][]{
    boxrule=0.5pt,
    breakable,
    colback=white,
    colbacktitle=white,
    colframe=black,
    fonttitle=\color{black},
    listing only,
    listing options={language=#3,frame=none,xleftmargin=0pt,title={#1}},
    sharp corners,
    title={Листинг~\thelstlisting~--~#1},
    title after break={Листинг~\thelstlisting\ (продолжение)},
    label={#2},
}

\newtcbinputlisting[use counter=lstlisting]{\flisting}[4][]{
    boxrule=0.5pt,
    breakable,
    colback=white,
    colbacktitle=white,
    colframe=black,
    fonttitle=\color{black},
    listing only,
    listing file={#3},
    listing options={language=#4,frame=none,xleftmargin=10pt},
    sharp corners,
    title={Листинг~\thelstlisting~--~#1},
    title after break={Листинг~\thelstlisting\ (продолжение)},
    label={#2}
}

% Подсчет общего количества рисунков, таблиц и источников
\newcounter{figures}
\regtotcounter{figures}
\AtBeginEnvironment{figure}{\stepcounter{figures}}

\newcounter{tables}
\regtotcounter{tables}
\AtBeginEnvironment{table}{\stepcounter{tables}}

\newcounter{references}
\regtotcounter{references}
\makeatletter
\pretocmd{\bibitem}{\stepcounter{references}}{}{}
\makeatother

\graphicspath{{figures/}{img/}}
